% Kamp III.2, Lemma 7: finite reductions across a fixed interior cut.

\subsubsection{Reducing a universal interior condition at a fixed cut}
\label{sec:p09-kamp-iii2-lemma7}

The next construction replaces a basic pattern reaching across an
interior point by finitely many pairs of patterns meeting at that point.
When one endpoint varies under a universal quantifier, a finite choice
of the fixed-side conditions determines the available disjuncts on the
varying side. Kamp's Lemma~7 uses this reduction with an additional
restriction: the varying-side disjunction contains no longer pattern
and at most one pattern as long as the original
\cite[ch.~III, lemma~7, pp.~80--83]{Kamp1968}.

\begin{hbconvention}{def:p09-kamp-cut-standing}{Patterns, lengths and finite aggregates}
Use the basic grammar and occurrence semantics of
\cref{def:p09-kamp-basic-classes,def:p09-kamp-basic-word}.
Fix a basic past pattern with occurrence word
\[
                W=(A_0,B_1,A_1,\ldots,B_n,A_n),\qquad n\geq1.
\]
Predicate letters may repeat. Each occurrence retains its point or
open-interval position. For any word with \(m\) intervals, put
\(\ell(W)=2m+1\); the length \(\ell(\phi)\) of a basic pattern
is the length of its occurrence word. Thus this length counts predicate
occurrences, rather than interval pieces or printed formula symbols.
The source makes this length convention explicit at the start of
Lemma~8, on p.~83.

For distinct endpoint names \(r,s\), write
\(\mathcal B_p(W;r,s)\) for the past-pattern recursion of
\cref{def:p09-kamp-labelled-intervals}, with the letters in \(W\)
as its unary labels. It is a member of \(\operatorname{Bp}_m(r,s)\).
All new bound names are fresh. In particular a pattern is obtained from
its occurrence word without changing its interval count.

For a finite family, an empty conjunction means the previously declared
\(T^K_w=P_0(w)\lor\neg P_0(w)\), and an empty disjunction means
\(F^K_w=P_0(w)\land\neg P_0(w)\), at the indicated variable.
These are formula abbreviations. An empty aggregate is not an additional
stage-zero basic pattern. Nonempty aggregates are associated from left
to right in their specified index order.
\end{hbconvention}

\begin{hbdefinition}{def:p09-kamp-cut-words}{All point cuts and gap cuts of one word}
Let \(\mathcal C_n\) be the finite index set
\[
 \mathcal C_n\defeq
 \{(\mathrm g,k):1\leq k\leq n\}\ \cup\
 \{(\mathrm p,k):1\leq k<n\}.
\]
Order it by increasing \(k\), with the gap tag before the point tag
at the same \(k\). At \(n=1\) it contains only \((\mathrm g,1)\).
There are \(2n-1\) entries. The tag \(\mathrm p\) means an
interior point position and \(\mathrm g\) an open-gap position;
these tags are not predicate letters or object-language constants.
For each entry \(c\), define two words \(W_c^L,W_c^R\).

For \(c=(\mathrm p,k)\), where \(1\leq k\leq n-1\), set
\[
\begin{aligned}
 W_c^L&=(A_0,B_1,A_1,\ldots,B_k,A_k),\\
 W_c^R&=(A_k,B_{k+1},A_{k+1},\ldots,B_n,A_n).
\end{aligned}
\]
The two words share the point label \(A_k\).
For \(c=(\mathrm g,k)\), where \(1\leq k\leq n\), set
\[
\begin{aligned}
 W_c^L&=(A_0,B_1,A_1,\ldots,B_{k-1},A_{k-1},B_k,B_k),\\
 W_c^R&=(B_k,B_k,A_k,B_{k+1},A_{k+1},\ldots,B_n,A_n).
\end{aligned}
\]
At \(k=1\), the first word is \((A_0,B_1,B_1)\); at \(k=n\),
the second is \((B_n,B_n,A_n)\). Each duplicate \(B_k\) has a
specific role: one occurrence labels the new endpoint and the other
labels its adjoining gap. This requires the old interval predicate at
the cut point itself.

Finally define the two basic formulas
\[
 L_c(r,s)\defeq\mathcal B_p(W_c^L;r,s),\qquad
 R_c(r,s)\defeq\mathcal B_p(W_c^R;r,s).
\]
Their displayed arguments identify their free endpoints; they are formula
abbreviations, not new relation symbols. Every word used here has at
least one interval. A point cut gives respectively \(k,n-k\)
intervals, and a gap cut gives \(k,n-k+1\).

These are the source's two word operations \(R_i^*,S_i^*\), written
with explicit occurrence positions. A split before an interior point
duplicates its point letter between the two words. A split before an
interval occurrence duplicates its interval letter inside each adjoining
word. Occurrence positions determine the case even when the same letter
also appears in the other kind of position
\cite[ch.~III, pp.~64, 81--82]{Kamp1968}.
\end{hbdefinition}

\begin{hblemma}{thm:p09-kamp-cut-factorization}{proved=annex:p9.kamp.cut.factorization}{Factorization at any strict interior point}
Let \(\phi\in\operatorname{Bp}_n(r,t)\) have word \(W\), and
let \(s\) be a distinct variable. Then
\[
 \mathsf{Lin}\models r<s<t\longrightarrow
 \left(\phi\longleftrightarrow
       \bigvee_{c\in\mathcal C_n}(L_c(r,s)\land R_c(s,t))\right).
\]
For each side, all listed basic patterns have length at most
\(2n+1\). Exactly one candidate on the left has that length,
namely \(L_{(\mathrm g,n)}\), and exactly one on the right has
that length, namely \(R_{(\mathrm g,1)}\). Every other candidate
on its respective side is shorter.

The equivalence holds on every nonempty strict linear order. There is
no density, completeness or endpointlessness assumption. In particular
the cut is an already supplied domain point, not a point created inside
an empty interval.
\end{hblemma}

\begin{hbmethod}{def:p09-kamp-cut-aggregates}{The finite alternatives for the two orientations}
Let \(x,y,u,v\) be distinct variable names, and let \(\psi_1,\psi_2\)
be arbitrary formulas of \(L\). Additional free parameters are allowed
and are left unchanged. The displayed outer \(\exists u\) and
inner \(\forall v\) have exactly their displayed scopes, so the
inner quantifier does not bind any free occurrence of \(v\) in
\(\psi_1\). Bound variables internal to newly constructed patterns
are chosen fresh from both formulas and from all four named variables.

List every subset \(J\subseteq\mathcal C_n\), including the empty
subset. Equivalently list all functions
\(\alpha:\mathcal C_n\to\{0,1\}\), with
\(J=\{c:\alpha(c)=1\}\). There are \(2^{2n-1}\) indexed
alternatives before any repetitions are removed.

For a pattern \(\phi(v,y)\) spanning the left interior interval,
form
\[
 \Sigma_J^L(u,y)\defeq\bigwedge_{c\in J}R_c(u,y),\qquad
 \rho_J^L(v,u)\defeq\bigvee_{c\in J}L_c(v,u).
\]
The empty cases are \(\Sigma_\varnothing^L=T^K_u\) and
\(\rho_\varnothing^L=F^K_v\).
For a pattern \(\phi(x,v)\) spanning the right interior interval,
form
\[
 \Sigma_J^R(x,u)\defeq\bigwedge_{c\in J}L_c(x,u),\qquad
 \rho_J^R(u,v)\defeq\bigvee_{c\in J}R_c(u,v).
\]
Here the empty cases are \(\Sigma_\varnothing^R=T^K_u\) and
\(\rho_\varnothing^R=F^K_v\).

Each conjunction concerns the fixed side of the cut. Each disjunction
concerns the endpoint that will vary under \(\forall v\).
The construction uses only the word \(W\), its two kinds of cuts,
finite conjunction and finite disjunction. In particular it does not
require a decision about the truth of \(\psi_1\) or \(\psi_2\).
Truth assignments to the fixed-side patterns are used to prove that the
complete list of alternatives works in every structure
\cite[ch.~III, pp.~81--83]{Kamp1968}.
\end{hbmethod}

\begin{hblemma}{thm:p09-kamp-iii2-lemma7-left}{proved=annex:p9.kamp.iii2.lemma7.left}{Kamp's left-interior reduction}
Suppose \(\phi\in\operatorname{Bp}_n(v,y)\), with \(n\geq1\),
and use its word to construct the left aggregates above. The formula
\[
 \exists u\left(x<u<y\land\psi_1\land
       \forall v\,(x<v<u\to\psi_2\lor\phi(v,y))\right)
\]
is equivalent over \(\mathsf{Lin}\) to
\[
 \bigvee_{J\subseteq\mathcal C_n}
 \exists u\left(x<u<y\land\psi_1\land\Sigma_J^L(u,y)\land
       \forall v\,(x<v<u\to\psi_2\lor\rho_J^L(v,u))\right).
\]
Every \(\Sigma_J^L\) is a finite conjunction of members of
\(\operatorname{Bp}(u,y)\), and every nonempty \(\rho_J^L\)
is a disjunction of members of \(\operatorname{Bp}(v,u)\).
Every such disjunct is also an elementary interval formula.
The empty aggregate conventions are as stated above.

Each \(\rho_J^L\) has at most one basic disjunct of length
\(\ell(\phi)\), and none of greater length. These bounds apply
to the actual disjunct list, including if different cuts happen to
produce equivalent formulas. No larger replacement bound is used
\cite[ch.~III, lemma~7(a), pp.~80--83]{Kamp1968}.
\end{hblemma}

\begin{hblemma}{thm:p09-kamp-iii2-lemma7-right}{proved=annex:p9.kamp.iii2.lemma7.right}{The right-interior reduction}
Suppose \(\phi\in\operatorname{Bp}_n(x,v)\), with \(n\geq1\),
and use its word to construct the right aggregates above. The formula
\[
 \exists u\left(x<u<y\land\psi_1\land
       \forall v\,(u<v<y\to\psi_2\lor\phi(x,v))\right)
\]
is equivalent over \(\mathsf{Lin}\) to
\[
 \bigvee_{J\subseteq\mathcal C_n}
 \exists u\left(x<u<y\land\psi_1\land\Sigma_J^R(x,u)\land
       \forall v\,(u<v<y\to\psi_2\lor\rho_J^R(u,v))\right).
\]
Each fixed-side condition \(\Sigma_J^R\) is a finite conjunction
of members of \(\operatorname{Bp}(x,u)\), and each nonempty
\(\rho_J^R\) is a finite disjunction of members of
\(\operatorname{Bp}(u,v)\), with the declared empty cases.
Each residual disjunction has at most one disjunct of length
\(\ell(\phi)\), and no disjunct of greater length.

This is the handbook's corrected aggregate formulation of the orientation
in Kamp's Lemma~7(b). The source refers its proof to the analogous left case;
the annex supplies both directions for this orientation in full
\cite[ch.~III, lemma~7(b), p.~81]{Kamp1968}.
\end{hblemma}

\begin{hbexample}{def:p09-kamp-cut-two-intervals}{All three cuts of a two-interval word}
For \(W=(A_0,B_1,A_1,B_2,A_2)\), the complete cut list is
\[
\begin{array}{c|c|c}
 c&W_c^L&W_c^R\\ \hline
 (\mathrm g,1)&(A_0,B_1,B_1)&(B_1,B_1,A_1,B_2,A_2)\\
 (\mathrm p,1)&(A_0,B_1,A_1)&(A_1,B_2,A_2)\\
 (\mathrm g,2)&(A_0,B_1,A_1,B_2,B_2)&(B_2,B_2,A_2).
\end{array}
\]
There are eight subset alternatives. Only the last left word has
length five, and only the first right word has length five.

To see why a single cut need not work for all universally quantified
points, take \(T=\{0,1,2,3,4,5,6\}\) in its usual order and set
\[
 A_0=\{1,2\},\quad B_1=\{3\},\quad A_1=\{2,4\},\quad
 B_2=\{3,4,5\},\quad A_2=\{6\}.
\]
These are the interpretations of five distinct predicate letters.
Let \(x=0,u=3,y=6\), take \(\psi_1=T^K_u\) and
\(\psi_2=F^K_v\). At \(v=1\), the realization \((1,2,6)\)
crosses the second gap; at \(v=2\), the realization \((2,4,6)\)
crosses the first. The point cut is unavailable because \(A_1(3)\)
is false. Neither gap cut alone works for both values of \(v\).
The alternative \(J=\{(\mathrm g,1),(\mathrm g,2)\}\) includes
both required residual disjuncts, and its two fixed-side conjuncts
both hold at \((3,6)\). The annex verifies every predicate test.
\end{hbexample}

\begin{hbexample}{def:p09-kamp-cut-empty-choice}{The empty alternative and the smallest word}
For \(W=(A,B,C)\), there is only the gap cut
\[
                    W^L=(A,B,B),\qquad W^R=(B,B,C).
\]
There are two alternatives, corresponding to the empty subset and
the singleton subset. Both words have the original length three.

If \(\psi_2\) holds throughout the universally tested interval,
the empty alternative applies regardless of the pattern letters.
For example, take \(0<1<2\), let every unary predicate be false,
and take \(\psi_1=T^K_u,\psi_2=T^K_v\). For \(x=0,y=2\)
the witness \(u=1\) makes each of the two original quantified
formulas true. The corresponding empty alternative has true fixed-side
condition and a false residual; the tautological \(\psi_2\)
supplies its universal condition. Every nonempty conjunction of basic
patterns is false in this structure because its initial point predicate
is false. Thus the empty alternative is substantive and cannot be
replaced by requiring a basic pattern on the fixed side.

An empty tested open interval likewise leaves the universal condition
vacuously true. The outer strict witness \(x<u<y\) still remains:
if \(x,y\) are adjacent, both displayed existential formulas and
every constructed alternative are false.
\end{hbexample}

\paragraph{Use in the following induction.}
The cut theorem and both universal reductions have the source's explicit
\(\mathsf{Lin}\) scope. Their restrictions to \(\mathsf{Com}\)
are therefore available for the following interval-property argument.
Their at-most-one bound controls the longest varying-side disjuncts
while retaining every shorter case. Lemma~8 begins on printed p.~83 and applies a
well-founded comparison of finite families to establish property
\(I(x,y)\) for a universal interior disjunction. That induction is
the next separate result; no conclusion of Lemma~8 is used here
\cite[ch.~III, pp.~83--84]{Kamp1968}.
